\chapter{Conclusions}
\label{chap:conclusions}

\section{Summary of Contributions}
This dissertation demonstrated the design and implementation of a C++20 exchange
infrastructure for benchmarking order book data structures. It explored the 
trade-offs to consider when designing order books as well as the implications 
of real-world system overheads. The key contributions of this work are as follows:

\begin{itemize}
    \item A modular C++20 exchange framework with interchangeable
    order book implementations and a TCP gateway for end-to-end benchmarking.
    \item A benchmarking suite that decomposes latency and
    throughput across direct and gateway modes, isolating data structure
    performance from system overheads.
    \item Empirical evidence on the performance trade-offs of direct indexing
    versus pointer-based order book strategies, and the limited impact of
    memory pooling on overall latency.
    \item A clear systems-level finding that in end-to-end gateway operation,
    network and transport overheads dominate latency, masking internal data
    structure differences that are visible in isolated direct mode.
\end{itemize}

These contributions advance our understanding of how order book design choices interact 
with system-level factors to influence performance in high frequency trading environments. 
The modular framework and comprehensive benchmarks provide a foundation for future 
research and practical exchange design decisions.

\section{Answers to Research Hypotheses}
In terms of hypotheses \ref{hyp:h1}--\ref{hyp:h3}, the evidence supports the
following conclusions. These conclusions are based on the repeated-run results
in Chapter~\ref{chap:testing_evaluation}, with confidence intervals and gateway
latency breakdown where relevant.

\paragraph{\ref{hyp:h1}: Direct indexing outperforms pointer-heavy alternatives in isolated mode.}
\textbf{Status: Supported.}

Direct-mode benchmarks show that the direct-indexed array outperforms alternative
structures across almost all scenarios, with mean latency around 210 ns.
Results also indicate that non-array structures remain clearly slower than
array even in shallow scenarios, reinforcing \ref{hyp:h1}.
This is reflected consistently in the direct-mode results
(Table~\ref{tab:throughput_results} and Table~\ref{tab:direct_mean_latency}).

\paragraph{\ref{hyp:h2}: Memory pooling reduces allocation jitter but does not remove structural complexity costs.}
\textbf{Status: Supported with Caveats.}

Results indicate that memory pooling improves performance in
selected scenarios, particularly \texttt{sparse\_extreme} and \texttt{mixed}.
At the same time, the pool logic introduces its own
structural cost in mixed workloads, where branch misprediction and match-path
complexity increase latency, and in \texttt{dense\_full} pool still trails map.
So pooling is clearly useful, but it is an
optimisation, not a full solution to data-structure complexity.
This trade-off is discussed directly in
Section~\ref{subsec:pool_hybrid_mech}.


\paragraph{\ref{hyp:h3}: In end-to-end gateway operation, system overhead masks internal data-structure gains.}
\textbf{Status: Supported.}
In gateway mode, network latency and queue latency dominate total
latency, while engine latency remains very small in comparison. For example, for the Array book the
network component is 2,014,045 ns out of 2,100,661 ns total, showing that most
end-to-end cost is outside the matching algorithm itself. This transition 
from a p < 0.0001 result to a p > 0.05 result proves that the measurement 
'noise' introduced by the system stack is statistically larger than the 
order book's 'signal.' This explains why
differences visible in direct mode are often compressed in gateway mode,
while still allowing scenario-specific exceptions (notably hybrid in
\texttt{dense\_full}) where queue/network accumulation diverges.
This pattern is captured in Table~\ref{tab:gateway_decomposition}.

\section{Critical Reflection}
This project provides strong evidence within a bounded experimental
scope, but should still be interpreted with care.

First, the benchmark setup captures the full software path from gateway to
engine, and the decomposition methodology gives visibility into
network latency, queue latency, and engine latency costs. However, the loopback environment does not include physical
network effects such as switch delay and NIC-level contention, so absolute
latency should be treated as a strict theoretical lower bound; in a physical 
production environment, the masking effects identified here would likely be even 
more pronounced due to hardware-level jitter.

Second, the scenario design is robust for comparing relative behaviour across
implementations, especially for highlighting dense-book stress effects. At the
same time, synthetic workloads cannot fully reproduce strategic behaviour seen
in live markets, so claims should remain bounded to the tested scenario family.

Third, the practical implication is that exchange engineers should avoid picking
data structures from microbenchmarks alone. Direct-mode results are necessary to
understand algorithmic efficiency, but gateway-mode results are necessary for
deployment decisions where system overhead dominates.

If this project were repeated, three design choices would be adjusted earlier.
First, while this study implemented rigorous pairwise Welch's t-tests to 
identify significant performance tiers, a more exhaustive evaluation plan 
would have included these formal tests for every single 
implementation-scenario combination from the first iteration. This would 
integrate statistical rigour even more deeply into the data-gathering phase.
Second, gateway experiments would include an additional host-level networking
baseline (for example, alternative socket or
transport configurations) to separate framework effects from OS-path effects.
Third, scenario generation would be extended with longer replay windows so that
steady-state and stress-state behaviour can be contrasted more explicitly.

\subsection{Methodological Reflection}
Reflecting on the engineering process, the decision to use a custom wait-free SPSC queue
was a successful trade-off for learning and performance. However, it introduced a sec-
ondary bottleneck during timing capture. In a future iteration, implementing a dedicated,
lock-free measurement buffer with deferred storage would minimise the "observer effect"
even further. Moreover, while the Hexagonal architecture provided excellent modular-
ity for swapping order books, the cost of virtual functions in the hot path is a known
architectural tax. Exploring CRTP (Curiously Recurring Template Pattern) for static
polymorphism could potentially reclaim several nanoseconds of matching-engine latency.
This is a key consideration for reaching the extreme low-latency threshold.

\section{Limitations}
\begin{itemize}
    \item Platform dependency: results were measured on a pinned WSL2/Zen4 host,
    so absolute numbers may change on other operating-system and CPU stacks.
    \item Workload realism: scenarios are representative and diverse, but still
    synthetic, and therefore cannot capture every market microstructure effect.
    \item Statistical scope: this dissertation uses confidence intervals and 
    pairwise Welch's t-tests for the main hypothesis checks, but does not 
    exhaustively test every implementation-scenario combination.
    \item Loopback boundary: gateway mode exercises software networking cost but
    excludes physical network path variability.
\end{itemize}

\section{Professional, Ethical, Social, and Legal Issues}
Developing sub-microsecond matching engines is as much an ethical challenge as a 
technical one. As this dissertation demonstrates, achieving O(1) matching via the 
ArrayOrderBook provides a massive competitive edge. This "race to zero" creates a 
significant "digital divide" in the markets; only large institutional players with the 
capital to implement hardware-aware optimisations and co-located infrastructure can 
realistically compete. Smaller firms and retail investors, who cannot afford isolated CPU 
cores and bespoke memory management, are left at a structural disadvantage \citep{lewis2014flash}.

There is also a clear conflict between raw speed and market stability. Stripping 
away branching logic and validation checks to reduce latency, as I explored in Chapter 4, 
introduces systemic risk. Algorithms that prioritise throughput over safety can propagate 
erroneous orders at nanosecond scales, potentially triggering "Flash Crashes" that human 
traders cannot stop in time \citep{kirilenko2017flash}. To counter this, the architecture 
built here enforces semantic parity and boundary checks—not just as good engineering, 
but as a necessary ethical safeguard.

Legally, modern exchange infra must sit within the MiFID II regulatory framework 
\citep{mifid2}. This requires "kill-switches" and robust testing prior to deploy-
ment. My benchmarking suite acts as a key compliance tool here; it allows engineers to 
measure the exact latency "tax" of these mandatory safety features, ensuring the engine 
remains stable and compliant without flying blind into production.

\section{Practical Recommendations}
Based on the empirical findings, the following recommendations are provided for exchange architects:

\begin{itemize}
    \item \textbf{Prioritise Determinism:} Use direct-indexed layouts in latency-critical paths, as they provide the most stable Matching-path performance under stress.
    \item \textbf{Targeted Optimisation:} Deploy memory pooling specifically for allocation-heavy scenarios (e.g., Sparse LOBs), as the structural overhead of pooling logic may degrade performance in denser scenarios.
    \item \textbf{System-Level Visibility:} Prioritise end-to-end latency decomposition. As shown in Chapter 5, internal engine gains are often masked by transport overhead, making networking the highest-impact area for further engineering investment.
\end{itemize}

\section{Future Work}
First, the 'masking effect' identified in gateway-mode benchmarks suggests 
that the next logical step is to move beyond the POSIX stack. Future research 
should integrate Kernel Bypass technologies, such as DPDK (Data Plane 
Development Kit) or user-space transport like OpenOnload. By eliminating 
the OS kernel from the hot path, researchers can evaluate whether the 
engine-level efficiencies identified in direct mode become visible at the 
end-to-end level.

Second, the stability of the ArrayOrderBook suggests it is a primary candidate 
for Hardware Acceleration. Implementing the price-map logic on an FPGA 
(Field Programmable Gate Array) would remove the CPU jitter and 
frequency-scaling penalties discussed in Chapter 5. A hybrid hardware-software 
architecture could offload the Price-Time priority matching to the FPGA while 
maintaining the FIX protocol logic in C++20, potentially reaching 
nanosecond-scale determinism across all market scenarios.

\section{Closing Statement}
Taken together, the objectives set in Chapter~\ref{chap:introduction} were met:
the framework enabled fair cross-implementation benchmarking, direct mode and
gateway mode exposed different bottlenecks, and the decomposition evidence
showed where end-to-end latency is actually spent. To ensure these findings 
were mathematically sound, every quantitative claim was verified through 
a formal statistical framework incorporating Welch's t-tests and 
Bonferroni corrections. This rigour demonstrates that the observed 
performance trends are scientifically robust and not attributable to 
measurement noise. The practical implication is
that data-structure engineering remains essential for engine efficiency, but
deployment-level latency improvements require equal attention to transport and
queueing overhead.

Overall, this dissertation demonstrates that order book performance is a 
system-level problem. In isolated mode, cache-friendly indexing gives clear 
wins, but in gateway operation, transport overhead can dominate and mask those 
gains. The framework and evidence produced here provide a practical basis for 
future research and exchange design decisions.
